\documentclass[10pt]{article}
\usepackage[margin=1in]{geometry}
\usepackage{booktabs}
\usepackage{amsmath}
\usepackage{url}
\usepackage[hidelinks]{hyperref}
\usepackage{longtable}
\setlength{\parindent}{0pt}
\setlength{\parskip}{5pt}

\title{Supplementary Material\\[2pt]
\large Ambient WiFi as a Radar Replacement for Automotive SLAM:\\
A Physics-Based Feasibility Study}
\author{Mulham Fetna \quad (ORCID 0009-0006-4432-798X)}
\date{}

\begin{document}
\maketitle

This supplement provides the artifact/reproducibility details and the full dataset
datasheet that support the main manuscript but exceed its page limit. It duplicates
no results; every number and figure in the paper is regenerated by the commands
below.

\section*{S1. Artifact availability and licensing}
All code, scene/run configurations, result artifacts, and the dataset are openly
available and permanently archived.
\begin{itemize}
\item \textbf{Repository:} \url{https://github.com/mulhamfetna/wifi-radar-slam}
\item \textbf{Zenodo concept DOI} (all versions): \texttt{10.5281/zenodo.21247288}
\item \textbf{License:} AGPL-3.0-or-later. \textbf{Environment:} Python $\geq$3.11;
Sionna~RT~2.0 (CPU/LLVM or CUDA) for simulation; scikit-learn for the discriminator;
CSIKit for the real-CSI proof-of-concept.
\item \textbf{Full quality:} prepend \texttt{WRS\_NUM\_SAMPLES=1000000} to any run.
\end{itemize}

\section*{S2. Reproducibility recipe}
Each element of the paper is produced by a single command from a released config:

\begin{longtable}{@{}p{5.4cm}p{9.2cm}@{}}
\toprule
Paper element & Command (from the repository root) \\
\midrule
Localization (Phase~A) & \texttt{python experiments/run\_phase\_a.py configs/nominal.yaml} \\
Operating envelope (Table~I) & \texttt{python experiments/run\_phase\_b.py} \\
Oracle map, controlled wall & \texttt{run\_phase\_a.py configs/controlled\_oracle.yaml} \\
Oracle map, reflective canyon & \texttt{run\_phase\_a.py configs/street\_metal\_oracle.yaml} \\
Realistic map (joint MUSIC) & \texttt{run\_phase\_a.py configs/controlled\_music\_joint.yaml} \\
60~GHz test (Table, Sec.~VII) & \texttt{run\_phase\_a.py configs/controlled\_music\_60ghz.yaml} and \texttt{...\_60ghz\_16ant.yaml} \\
Scene render (Fig.~2) & \texttt{python experiments/render\_scene.py} \\
Dataset (WiFiSLAM-Sim) & \texttt{python experiments/make\_dataset.py configs/nominal.yaml data/wifislam\_sim.npz} \\
Learned discriminator & \texttt{python experiments/train\_discriminator.py data/wifislam\_sim.npz} \\
Real-CSI proof-of-concept & \texttt{sh experiments/fetch\_real\_csi.sh} then \texttt{python experiments/run\_real\_csi.py} \\
\bottomrule
\end{longtable}

\section*{S3. WiFiSLAM-Sim dataset datasheet}
A ray-traced outdoor/vehicular WiFi-CSI dataset (no public outdoor/vehicular
WiFi-CSI dataset otherwise exists). \emph{Simulated} (Sionna~RT), not hardware;
intended for algorithm development and as the labelled substrate for path
discrimination. Generated by \texttt{make\_dataset.py}; loaded with
\texttt{wifi\_radar\_slam.dataset.CsiDataset.load}.

\begin{tabular}{@{}lll@{}}
\toprule
key & shape & description \\
\midrule
\texttt{csi} & $(n_f, n_{ap}, n_{rx}, n_{sc})$ complex64 & noisy channel frequency response \\
\texttt{poses} & $(n_f, 3)$ & ground-truth vehicle $x,y,\text{yaw}$ (m, m, rad) \\
\texttt{ap\_positions} & $(n_{ap}, 3)$ & access-point coordinates (m) \\
\texttt{gt\_map} & $(M, 2)$ & ground-truth facade footprints (m) \\
\texttt{paths} & $(P, 9)$ & flat oracle path table (below) \\
\texttt{meta} & JSON & carrier/bandwidth/subcarriers/antennas/scene/SNR/units \\
\bottomrule
\end{tabular}

\textbf{\texttt{paths} columns:} \texttt{[frame, ap, delay\_s, phi\_r, theta\_r,
n\_bounce, first\_interaction\_type, object\_id, is\_floor]}. The last three are the
ground-truth \emph{labels} for learned LOS/reflection/floor discrimination
(\texttt{first\_interaction\_type}: 0 none, 1 specular, 2 diffuse, 4 refraction,
8 diffraction). \textbf{Tasks:} localization (\texttt{csi}$\to$\texttt{poses}),
mapping (reflectors vs \texttt{gt\_map}), path discrimination (features from
\texttt{paths}). \textbf{Limitations:} simulated (AWGN only); one straight trajectory
per record; sub-7~GHz unless generated from a 60~GHz config.

\section*{S4. Real commodity-CSI provenance}
The front-end proof-of-concept uses real captures shipped with CSIKit
(\url{https://github.com/Gi-z/CSIKit}, under its licence): Intel~5300
(\texttt{log.all\_csi.6.7.6.dat}; 30 subcarriers, 3 antennas) and Broadcom nexmon
802.11ac (\texttt{example\_4358.pcap}). Fetch them with
\texttt{experiments/fetch\_real\_csi.sh}. These are indoor/static captures with no
ground-truth trajectory; they validate that the sensing front-end runs on measured
commodity CSI, not full-system SLAM.

\end{document}
